\wegmannvsstyledistancetable
\section{Related Work}
\label{sec:relatedwork}

\paragraph{Style Representations} In previous work, style representations were learned from unlabeled texts \citep[inter alia]{styletransfervector2,styletransfervector6,textsettr}. Due to the lack of parallel datasets covering diverse style features, \citet{deepstyle} and \citet{styleemb} have both recently used authorship as a proxy for style as we discussed previously. \citet{luar} trained
embeddings to uniquely represent different authors. Although these embeddings capture  features representative of authors' style, they also capture content features due to the lack of control for content-related aspects. \citet{lisa}  trained LISA, a style vector created by annotating texts with their stylistic features using LLMs. LISA, however, trades off performance to create a vector with interpretable dimensions. The highest quality style representations to date come from approaches that employ a contrastive learning objective \citep{styleemb}. In this paper, we propose leveraging the strengths of generative LLMs to build a dataset that will serve to train strong content-independent style embeddings using a contrastive learning objective.  From a practical perspective, style representations are useful in downstream applications such as arbitrary text style transfer \citep{learningtogenerate,paraguide,tinystyler} where they help steer and guide transfer, and content-independence is important for reducing hallucinations in generations.

\paragraph{LLMs and Text Style} LLMs are commonly used for style transfer \citep{reif,promptandrerank,lowresourceauthorshipstyletransfer,styleremix} and style analysis \citep{iclef,lisa}. Their strength with style-related tasks has been demonstrated, yet leveraging this knowledge to build strong content-independent style representations has not yet been explored.
